\subsection{Determinante von $n\times n$-Matrizen}
\label{subsec:linear_systems_laplace_expansion}

Bei drei Unbekannten tritt eine $3\times3$-Determinante auf. Wir führen ihre
Berechnung zunächst konkret an der ersten Zeile durch. Zu jedem Eintrag wird
seine Zeile und seine Spalte gestrichen; die vier verbleibenden Einträge bilden
eine $2\times2$-Determinante.

\begin{center}
  \CachedTikzFigure{content/resources/chapters/linear_systems/tikz/three_by_three_laplace_expansion.tex}
\end{center}

Für einen Eintrag $a_{ij}$ wird seine Zeile $i$ und seine Spalte $j$ entfernt.
Die verbleibende $(n-1)\times(n-1)$-Matrix bezeichnen wir mit $A_{ij}$.

\begin{TheoremBox}{Laplacescher Entwicklungssatz}{linear_systems_laplace_expansion}
  Eine $n\times n$-Determinante kann nach jeder beliebigen Zeile oder Spalte
  entwickelt werden. Dabei gehen wir schrittweise vor:
  \begin{enumerate}[label=\textbf{Schritt \arabic*:},leftmargin=7.2em,itemsep=0.7em]
    \item Die gedachten Vorzeichen werden nach dem Schachbrettmuster
      eingetragen:
      \[
        \begin{pmatrix}
          +&-&+&\cdots\\
          -&+&-&\cdots\\
          +&-&+&\cdots\\
          \vdots&\vdots&\vdots&\ddots
        \end{pmatrix}.
      \]
    \item Eine beliebige Zeile oder Spalte wird für die Entwicklung ausgewählt.
    \item Entlang dieser Zeile oder Spalte wird für jeden Eintrag $a_{ij}$ der
      Summand
      \[
        \text{Vorzeichen an der Stelle }(i,j)
        \;\cdot\;a_{ij}\;\cdot\;\det(A_{ij})
      \]
      gebildet. Die Summe aller Summanden ist $\det(A)$.
  \end{enumerate}
  Dieses Vorgehen heißt \emph{Laplace-Entwicklung}
  \EnglishTerm{Laplace expansion}.
\end{TheoremBox}

\begin{ExampleBox}{Eine $3\times3$-Determinante nach Zeile und Spalte entwickeln}{linear_systems_three_by_three_determinant}
  Wir berechnen dieselbe Determinante auf zwei Wegen.
  \begin{enumerate}[label=\alph*),leftmargin=1.8em,itemsep=0.8em]
    \item \textbf{Entwicklung nach der ersten Zeile:}
      \[
        D=\det\!\begin{pmatrix}
          \IfClozeTF{1}{\colorbox{yellow!30}{\(\mathstrut 1\)}}&
          \IfClozeTF{2}{\colorbox{yellow!30}{\(\mathstrut 2\)}}&
          \IfClozeTF{3}{\colorbox{yellow!30}{\(\mathstrut 3\)}}\\
          4&5&6\\
          7&8&0
        \end{pmatrix}.
      \]
      \begin{ClozeArea}[32]{12}
        \[
          \begin{aligned}
            D
            &={+}\,1\det\!\begin{pmatrix}5&6\\8&0\end{pmatrix}
              -2\det\!\begin{pmatrix}4&6\\7&0\end{pmatrix}
              +3\det\!\begin{pmatrix}4&5\\7&8\end{pmatrix}\\[0.7em]
            &=1(5\cdot0-6\cdot8)
              -2(4\cdot0-6\cdot7)
              +3(4\cdot8-5\cdot7)\\[0.7em]
            &=-48+84-9=27.
          \end{aligned}
        \]
      \end{ClozeArea}
    \item \textbf{Entwicklung nach der zweiten Spalte:}
      \[
        D=\det\!\begin{pmatrix}
          1&\IfClozeTF{2}{\colorbox{yellow!30}{\(\mathstrut 2\)}}&3\\
          4&\IfClozeTF{5}{\colorbox{yellow!30}{\(\mathstrut 5\)}}&6\\
          7&\IfClozeTF{8}{\colorbox{yellow!30}{\(\mathstrut 8\)}}&0
        \end{pmatrix}.
      \]
      \begin{ClozeArea}[32]{12}
        \[
          \begin{aligned}
            D
            &=-2\det\!\begin{pmatrix}4&6\\7&0\end{pmatrix}
              +5\det\!\begin{pmatrix}1&3\\7&0\end{pmatrix}
              -8\det\!\begin{pmatrix}1&3\\4&6\end{pmatrix}\\[0.7em]
            &=-2(4\cdot0-6\cdot7)
              +5(1\cdot0-3\cdot7)
              -8(1\cdot6-3\cdot4)\\[0.7em]
            &=84-105+48=27.
          \end{aligned}
        \]
      \end{ClozeArea}
  \end{enumerate}
\end{ExampleBox}

\begin{KeyPointBox}{Günstig entwickeln}{linear_systems_laplace_choose_row_or_column}
  Das Ergebnis ist unabhängig von der gewählten Zeile oder Spalte. Zweckmäßig
  ist eine Zeile oder Spalte mit möglichst vielen Nullen: Jeder Eintrag null
  lässt einen vollständigen Summanden entfallen.
\end{KeyPointBox}

\begin{ExampleBox}{Nullen gezielt nutzen}{linear_systems_laplace_sparse_column}
  In der dritten Spalte ist nur ein Eintrag von null verschieden. Daher
  entwickeln wir zuerst nach dieser Spalte und anschließend nach der ersten
  Spalte der verbleibenden $3\times3$-Determinante:
  \[
    D=\det\!\begin{pmatrix}
      1&2& 0&4\\
      0&7& 0&2\\
      0&5& 0&4\\
      3&3&-1&2
    \end{pmatrix}.
  \]
  \begin{ClozeArea}[30]{13}
    \[
      \begin{aligned}
        D
        &=(-1)^{4+3}(-1)
          \det\!\begin{pmatrix}
            1&2&4\\
            0&7&2\\
            0&5&4
          \end{pmatrix}\\[0.8em]
        &=1\det\!\begin{pmatrix}7&2\\5&4\end{pmatrix}
          =7\cdot4-2\cdot5
          =18.
      \end{aligned}
    \]
  \end{ClozeArea}
\end{ExampleBox}
